\chapter{Repérage, distances et lieux géométriques}
\section*{Exercice 7.1 -- Distance et milieu}
$\overrightarrow{AB}$ a pour différences de coordonnées $(6,8)$, donc $AB=\sqrt{36+64}=10$. Le milieu est $I(2,6)$, obtenu en faisant la moyenne de chaque paire de coordonnées.

\section*{Exercice 7.2 -- Triangle rectangle}
$AB=4$, $AC=3$, $BC=\sqrt{16+9}=5$. Comme $AB^2+AC^2=BC^2$, la réciproque de Pythagore donne un angle droit en $A$. Cela se lit aussi sur les axes perpendiculaires.

\section*{Exercice 7.3 -- Cercle}
Le centre est $(2,-1)$ et le rayon $5$. Pour $(5,3)$, la somme vaut $3^2+4^2=25$, donc le point est sur le cercle.

\section*{Exercice 7.4 -- Médiatrice}
L'égalité des carrés des distances donne $(x+2)^2+y^2=(x-4)^2+y^2$, soit $12x=12$. Le lieu est $x=1$, droite verticale passant par le milieu de $[AB]$, donc sa médiatrice.

\section*{Exercice 7.5 -- Distance à un axe}
Les distances sont $4$ à l'axe horizontal et $3$ au vertical. Pour $P(t,0)$, $MP^2=(t-3)^2+16\geq16$, avec égalité uniquement si $t=3$.

\section*{Exercice 7.6 -- Rectangle général}
Prendre $(0,0),(a,0),(a,b),(0,b)$ dans un repère orthonormé. Les deux milieux valent $(a/2,b/2)$ ; chaque diagonale a pour carré de longueur $a^2+b^2$. On a traité tous les rectangles avec $a,b>0$.
